\chapter{Tastenkürzel und Mausbedienung}\label{kap:tastenkuerzel}
\sloppy % lange Befehls-Aufforderungen in Schreibmaschinenschrift umbrechen
\index{Tastenkürzel}

Dieser Anhang fasst alle Tastenkürzel, Maus- und Trackpad-Gesten von \produkt{} zusammen,
gegliedert nach Bereichen. Die Angaben sind dem Programmcode entnommen. Wo eine Funktion im
Handbuch ausführlich beschrieben ist, verweist die Tabelle auf das Kapitel.

\section{Strg und Cmd, Alt und Option}\label{sec:strg-cmd}
\index{Cmd-Taste}\index{Strg-Taste}

Das Handbuch schreibt Tastenkürzel in der Form von Windows und Linux, also mit \keys{Strg}.
Unter \textbf{macOS} verwenden Sie statt \keys{Strg} die \textbf{Befehlstaste} (Cmd): Aus
\keys{Strg+S} wird \keys{Cmd+S}, aus \keys{Strg}+\maus{Klick} wird \keys{Cmd}+\maus{Klick}. Das
gilt für Menükürzel ebenso wie für die Auswahl mit der Maus und das Zoomen mit
\keys{Strg}+Scrollen. Die Taste \keys{Alt} heißt auf dem Mac \keys{Option}; \keys{Shift} ist die
Umschalttaste. Die Rücktaste (auf dem Mac mit \enquote{delete} beschriftet) löscht wie
\keys{Entf} die Auswahl.

Die Menüs zeigen neben jedem Eintrag das auf Ihrem System gültige Kürzel an. Die Einträge, die
in den Tabellen mit \enquote{Systemstandard} markiert sind, verwenden die Standardkürzel des
Betriebssystems; sie können sich zwischen Windows, macOS und Linux geringfügig unterscheiden --
maßgeblich ist dann die Anzeige im Menü.

\begin{hinweis}
Einbuchstabige Kürzel wie \keys{H} oder \keys{Z}, \keys{E} wirken, solange eine Ansicht (oder
ein anderes Element des Hauptfensters außer einem Textfeld) den Tastaturfokus hat. Steht der
Cursor in der Befehlszeile oder in einem Eingabefeld, werden die Buchstaben ganz normal
eingegeben. \keys{H} in einer Ansicht blendet nur aus, wenn etwas ausgewählt ist; ohne Auswahl
wird der Buchstabe in die Befehlszeile geschrieben (etwa für den Alias \code{h} der Helix).
\end{hinweis}

% ---------------------------------------------------------------------------
\section{Datei}\label{sec:tk-datei}

\begin{referenz}{>{\raggedright\arraybackslash}p{0.2\linewidth}>{\raggedright\arraybackslash}p{0.24\linewidth}>{\raggedright\arraybackslash}p{0.16\linewidth}X}
\toprule
Kürzel & Menü & Bemerkung & Wirkung \\
\midrule\endhead
\keys{Strg+N} & \menu{File,New} & Systemstandard & Neues, leeres Dokument. \\
\keys{Strg+O} & \menu{File,Open…} & Systemstandard & Dokument oder STEP-Datei öffnen. \\
\keys{Strg+I} & \menu{File,Import STEP…} & & STEP-Datei zum Dokument hinzufügen. \\
\keys{Strg+S} & \menu{File,Save} & Systemstandard & Dokument speichern. \\
\keys{Strg+Shift+S} & \menu{File,Save As…} & Systemstandard & Unter neuem Namen speichern. \\
\keys{Strg+Q} & \menu{File,Quit} & Systemstandard & Programm beenden (auf dem Mac im
  Programmmenü). \\
-- & \menu{File,Export STEP…} & & ohne Kürzel \\
-- & \menu{File,Export STL…} & & ohne Kürzel \\
\bottomrule
\end{referenz}

Einzelheiten zu den Dateiformaten: Kapitel~\ref{kap:dateien}.

% ---------------------------------------------------------------------------
\section{Bearbeiten}\label{sec:tk-bearbeiten}

\begin{referenz}{l>{\raggedright\arraybackslash}p{0.27\linewidth}X}
\toprule
Kürzel & Menü & Wirkung \\
\midrule\endhead
\keys{Strg+Z} & \menu{Edit,Undo} & Letzten Schritt rückgängig machen (Systemstandard). \\
\keys{Strg+Shift+Z} & \menu{Edit,Redo} & Rückgängig gemachten Schritt wiederholen
  (Systemstandard; unter Windows ist üblicherweise auch \keys{Strg+Y} belegt). \\
\keys{Entf}, Rücktaste & \menu{Edit,Delete} & Ausgewählte Objekte und Ergebnis-Zeichnungen
  (Strahlen, Kamera) löschen (rückgängig machbar). Einzelne Flächen, Kanten und Ecken lassen sich
  nicht löschen. \\
\keys{Strg+A} & \menu{Edit,Select All} & Alle sichtbaren Objekte auswählen. \\
\keys{Strg+G} & \menu{Edit,Group…} & Auswahl gruppieren (Kapitel~\ref{kap:organisieren}). \\
\keys{Strg+Shift+G} & \menu{Edit,Ungroup} & Gruppe auflösen. \\
\keys{H} & \menu{Edit,Hide} & Ausgewählte Objekte ausblenden (auch Knopf in der
  Ansichts-Werkzeugleiste). Ohne Auswahl geht \keys{H} in einer Ansicht in die Befehlszeile. \\
-- & \menu{Edit,Isolate selected} & Nur die Auswahl anzeigen; ohne Kürzel. \\
\keys{Shift+H} & \menu{Edit,Unhide} & Mit Hide oder Isolate ausgeblendete Objekte wieder
  anzeigen. Layer und Kontrollkästchen im Modellbaum behalten ihren Zustand. \\
-- & \menu{Edit,Show all} & Alles anzeigen, auch ausgeschaltete Layer und Modellbaum-Einträge;
  ohne Kürzel. \\
\keys{Esc} & \menu{Edit,Select None / Cancel} & Der Reihe nach: laufenden Befehl abbrechen,
  sonst laufenden Plugin-Auftrag abbrechen, sonst Auswahl aufheben. \\
\bottomrule
\end{referenz}

% ---------------------------------------------------------------------------
\section{Ansicht und Navigation}\label{sec:tk-ansicht}
\index{Navigation}

\subsection{Tasten}

\begin{referenz}{>{\raggedright\arraybackslash}p{0.38\linewidth}X}
\toprule
Kürzel & Wirkung \\
\midrule\endhead
\keys{Z}, dann \keys{E} & \gui{Zoom extents}: das ganze Modell in der aktiven Ansicht
  einpassen. \\
\keys{Z}, dann \keys{S} & \gui{Zoom selected}: die ausgewählten Elemente füllen die aktive
  Ansicht; ohne Auswahl wie Zoom extents. \\
\keys{Leertaste} & aktive Ansicht maximieren bzw.\ wieder das $2 \times 2$-Raster zeigen. \\
\keys{Shift}+\maus{Klick} auf \gui{Wireframe}, \gui{Surfaces}, \gui{Shaded}, \gui{Zoom
  extents} oder \gui{Zoom selected} (Werkzeugleiste) & wirkt auf alle vier Ansichten statt nur
  auf die aktive. \\
\bottomrule
\end{referenz}

Die Darstellungsmodi und die Docks haben keine Tastenkürzel; Sie erreichen sie über die
Werkzeugleiste und das Menü \menu{View}.

\subsection{Maus}

\begin{referenz}{p{0.3\linewidth}p{0.3\linewidth}X}
\toprule
Aktion & Ansicht Perspective & Ansichten Front, Left, Top \\
\midrule\endhead
Drehen (Orbit) & \maus{rechts ziehen} oder \keys{Alt}+\maus{links ziehen} & -- (kein Drehen) \\
Verschieben (Pan) & \maus{mittlere Taste ziehen} oder \keys{Shift}+\maus{rechts ziehen} &
  \maus{mittlere Taste ziehen}, \maus{rechts ziehen} oder \keys{Alt}+\maus{links ziehen} \\
Zoomen zum Cursor & \maus{Mausrad} & \maus{Mausrad} \\
Zoom extents & \maus{Doppelklick mit der mittleren Taste} & ebenso \\
Ansichtsmenü & \maus{Klick auf den Ansichtstitel} oben links & ebenso \\
\bottomrule
\end{referenz}

Das Ansichtsmenü bietet die Standardansichten Front, Left, Top und Perspective, den Schalter
\gui{Perspective projection}, \gui{Zoom extents}, \gui{Zoom selected} und
\gui{Maximize / restore  (Space)}. Die Ansicht, in die Sie zuletzt geklickt haben, ist die
\emph{aktive} Ansicht. Ausführlich: Kapitel~\ref{kap:navigation}.

\subsection{Trackpad}

\begin{referenz}{>{\raggedright\arraybackslash}p{0.38\linewidth}X}
\toprule
Geste & Wirkung \\
\midrule\endhead
\maus{Zwei-Finger-Scrollen} & Ansicht verschieben (Pan). \\
\keys{Strg}+\maus{Zwei-Finger-Scrollen} & zoomen zum Cursor. \\
\maus{Aufziehen/Zusammenziehen} (Pinch) & zoomen zum Cursor. \\
\maus{Doppeltippen mit zwei Fingern} & Zoom extents (Smart Zoom, macOS). \\
\bottomrule
\end{referenz}

% ---------------------------------------------------------------------------
\section{Auswahl}\label{sec:tk-auswahl}
\index{Auswahl}

Diese Bedienung gilt, solange kein Befehl einen Punkt erwartet. Während der Punkteingabe setzt
ein Linksklick stattdessen einen Punkt.

\begin{referenz}{>{\raggedright\arraybackslash}p{0.38\linewidth}X}
\toprule
Eingabe & Wirkung \\
\midrule\endhead
\maus{Klick} auf ein Objekt & Objekt auswählen (ersetzt die Auswahl). Ein Klick auf eine
  Ergebnis-Zeichnung (Strahlen eines Laufs, Kamera des Path Tracers) wählt sie als Ganzes aus. \\
\maus{Klick} ins Leere & Auswahl aufheben. \\
\keys{Shift}+\maus{Klick} & Objekt zur Auswahl hinzufügen. \\
\keys{Strg}+\maus{Klick} & Objekt aus der Auswahl entfernen. \\
\keys{Strg+Shift}+\maus{Klick} & einzelne Fläche, Kante oder Ecke eines Objekts umschalten
  (Teilelement-Auswahl). Mit gehaltenem \keys{Strg+Shift} hebt die Hervorhebung unter dem
  Cursor das Teilelement hervor. \\
\maus{links ziehen} von links nach rechts & Fensterauswahl: Objekte, die ganz im Rechteck
  liegen. \\
\maus{links ziehen} von rechts nach links & Kreuzungsauswahl: auch Objekte, die das Rechteck
  berühren. \\
\keys{Shift} / \keys{Strg} beim Rechteck & Objekte hinzufügen bzw.\ entfernen. \\
\keys{Strg+A} & alle sichtbaren Objekte auswählen. \\
\keys{Esc} & Auswahl aufheben (wenn weder ein Befehl noch ein Plugin-Auftrag läuft). \\
\bottomrule
\end{referenz}

Die Reihenfolge, in der Sie Objekte auswählen, ist für einige Befehle wichtig (z.\,B.\
Boolean difference, Loft, Sweep, Trim); siehe Abschnitt~\ref{sec:auswahlbefehle}. Im Modellbaum
gelten die üblichen Listenregeln: \keys{Shift}+\maus{Klick} wählt einen Bereich,
\keys{Strg}+\maus{Klick} einzelne Einträge.

% ---------------------------------------------------------------------------
\section{Befehle und Befehlszeile}\label{sec:tk-befehle}

\begin{referenz}{>{\raggedright\arraybackslash}p{0.38\linewidth}X}
\toprule
Eingabe & Wirkung \\
\midrule\endhead
Tippen in einer Ansicht & Die Zeichen gehen in die Befehlszeile (Befehlsnamen, Koordinaten,
  Zahlen). \\
\keys{Enter} & Eingabe der Befehlszeile ausführen; bei leerer Zeile den laufenden Befehl
  beenden oder, wenn keiner läuft, den letzten Befehl wiederholen. \\
\maus{Rechtsklick} (ohne Ziehen) in einer Ansicht & laufenden Befehl beenden (wie
  \keys{Enter}); läuft keiner, den letzten Befehl wiederholen. Rechts \emph{ziehen} dreht bzw.\
  verschiebt nur die Ansicht. \\
\keys{Esc} & laufenden Befehl abbrechen. \\
Alias + \keys{Enter} & Befehl starten, z.\,B.\ \code{l} Line, \code{c} Circle (centre,
  radius), \code{m} Move, \code{h} Helix. Der Tooltip jedes Werkzeugs nennt den Alias
  (\enquote{Type: l + Enter}). \\
\keys{Entf} & Werkzeug Delete: Auswahl löschen. \\
\keys{F8} & Ortho ein/aus (Knopf \gui{Ortho} in der Statusleiste). \\
\maus{Klick} auf einen Werkzeugknopf & Befehl starten. \\
\maus{Klick} auf eine Gruppenüberschrift der Werkzeugspalte & Gruppe auf- oder zuklappen. \\
\bottomrule
\end{referenz}

Die Befehle haben keine eigenen Tastenkürzel; ein einzelner Buchstabe wie \keys{L} startet
nichts. Er wird -- wie jedes in einer Ansicht getippte Zeichen -- in die Befehlszeile
geschrieben und erst mit \keys{Enter} als Alias ausgeführt. Einziges Werkzeug mit echtem
Tastenkürzel ist \gui{Delete} (\keys{Entf}). Alle Befehle mit Namen und Aliassen stehen in
Anhang~\ref{kap:befehle}; die Eingabe von Punkten und Zahlen beschreibt
Abschnitt~\ref{sec:punkteingabe}.

% ---------------------------------------------------------------------------
\section{Objektfang}\label{sec:tk-fang}

Diese Eingaben wirken, während ein Befehl einen Punkt erwartet (Abschnitt~\ref{sec:objektfang}).

\begin{referenz}{>{\raggedright\arraybackslash}p{0.38\linewidth}X}
\toprule
Eingabe & Wirkung \\
\midrule\endhead
\keys{F3} & Objektfang insgesamt ein/aus (Knopf \gui{Snap} in der Statusleiste). \\
\keys{Tab} & nächster Fangkandidat unter dem Cursor. \\
\keys{Shift+Tab} & vorheriger Fangkandidat. \\
\keys{Alt} gedrückt halten & Fang unterdrücken. \\
\keys{Alt}+\maus{Klick} (ohne Ziehen) & freien Punkt ohne Fang setzen. \\
\keys{Alt}+\maus{links ziehen} & Ansicht drehen (Perspective) bzw.\ verschieben, auch während
  der Punkteingabe. \\
\maus{Klick} auf \gui{Point}, \gui{End}, \gui{Corner}, \gui{Int}, \gui{Cen}, \gui{Quad},
  \gui{Mid}, \gui{Edge}, \gui{Near}, \gui{Grid} & einzelne Fangart ein/aus (Statusleiste). \\
\bottomrule
\end{referenz}

% ---------------------------------------------------------------------------
\section{Node-Editor}\label{sec:tk-node-editor}
\index{Node-Editor!Bedienung}

Der Node-Editor (Dock \gui{Node workflow}, Kapitel~\ref{kap:workflow}) wird wie die
CAD-Ansichten bedient. Die Tasten wirken, solange der Editor den Tastaturfokus hat; \keys{Z},
\keys{E}, \keys{Z}, \keys{S} und \keys{Pos1} gehören dann dem Editor und nicht den CAD-Ansichten.

\begin{referenz}{>{\raggedright\arraybackslash}p{0.38\linewidth}X}
\toprule
Eingabe & Wirkung \\
\midrule\endhead
\keys{Z}, dann \keys{E} oder \keys{Pos1} & alle Blöcke einpassen. \\
\keys{Z}, dann \keys{S} & ausgewählte Blöcke einpassen (ohne Auswahl: alle). \\
\keys{Esc} & eine Ebene aus einer geöffneten Gruppe zurück; während der Aufnahme eines Zweigs
  (\gui{New branch}): Aufnahme abbrechen. Sonst wirkt \keys{Esc} wie im Hauptfenster. \\
\maus{Mausrad} & zoomen zum Cursor. \\
\keys{Strg}+\maus{Scrollen}, \maus{Pinch} & zoomen zum Cursor (Trackpad). \\
\maus{Zwei-Finger-Scrollen} & verschieben (Trackpad). \\
\maus{Doppeltippen mit zwei Fingern} & alle Blöcke einpassen. \\
\maus{mittlere Taste ziehen}, \maus{rechts ziehen} oder \keys{Alt}+\maus{links ziehen} &
  verschieben. \\
\maus{Doppelklick mit der mittleren Taste} & alle Blöcke einpassen. \\
\maus{Rechtsklick} (ohne Ziehen) & Kontextmenü des Blocks, der Verbindung oder der freien
  Fläche unter dem Cursor (auch mit der Kontextmenü-Taste der Tastatur). \\
\maus{Klick} auf einen Block & Block auswählen; \maus{Ziehen} verschiebt ihn. \\
\maus{links ziehen} auf freier Fläche & Auswahlrechteck für Blöcke. \\
\maus{Ziehen} von einem Anschluss zu einem anderen & Verbindung herstellen. \\
\maus{Doppelklick} auf einen Block & Anzeige-Block: sein Fenster öffnen; Gruppe: hinein
  wechseln; sonst: Parameter bearbeiten. \\
\maus{Klick} auf den Ebenennamen oben links & eine Ebene aus der Gruppe zurück. \\
\bottomrule
\end{referenz}

Blöcke, Verbindungen und Varianten löschen Sie über das Kontextmenü (\gui{Delete block},
\gui{Delete link}, \gui{Delete connection}). Die Taste \keys{Entf} hat im Node-Editor keine
eigene Funktion; sie wirkt auf die Auswahl im CAD-Modell (\menu{Edit,Delete}).

In der Zweigliste neben dem Editor öffnet ein \maus{Doppelklick} auf einen Schritt die Auswahl
seines Kanals.

% ---------------------------------------------------------------------------
\section{Bildfenster}\label{sec:tk-bild}
\index{Bildfenster}

Für das Fenster \gui{Image} und die Fenster der Anzeige-Blöcke mit Bildern:

\begin{referenz}{>{\raggedright\arraybackslash}p{0.38\linewidth}X}
\toprule
Eingabe & Wirkung \\
\midrule\endhead
\maus{Mausrad} & zoomen zum Cursor (höchstens bis auf wenige Bildpunkte, nicht über das ganze
  Bild hinaus). \\
\maus{links ziehen} & verschieben. \\
\maus{Doppelklick} & ganzes Bild zeigen. \\
Maus über dem Bild & Koordinaten und Wert des Bildpunkts unter dem Cursor werden angezeigt. \\
Kontrollkästchen \gui{log} & logarithmische Farbskala. \\
\bottomrule
\end{referenz}

% ---------------------------------------------------------------------------
\section{Plot und Tabelle}\label{sec:tk-plot}

Die Fenster \gui{Plot} und \gui{Table} sowie die Plot-Fenster der Anzeige-Blöcke sind reine
Anzeigen: Sie haben keine Tastenkürzel und keine Mausgesten zum Zoomen oder Verschieben. Die
Tabelle ist nicht editierbar. Welche Kurven sichtbar sind, steuern Sie im Dock \gui{Results}
(Kapitel~\ref{kap:plugins}).

% ---------------------------------------------------------------------------
\section{Weitere Mausbedienung in den Docks}\label{sec:tk-docks}

\begin{referenz}{>{\raggedright\arraybackslash}p{0.38\linewidth}X}
\toprule
Eingabe & Wirkung \\
\midrule\endhead
\maus{Doppelklick} auf die Spalte \gui{Colour} im Dock \gui{Layers} & Farbe des Layers
  wählen. \\
\maus{Rechtsklick} im Dock \gui{Model} & Kontextmenü: \gui{Group selection…}, \gui{Ungroup},
  \gui{Hide}, \gui{Isolate}, \gui{Unhide}, \gui{Show all}. \\
\maus{Rechtsklick} im Dock \gui{Results} & Kontextmenü der Ergebnisse. \\
\maus{Doppelklick} auf einen Listeneintrag in einem Plugin-Formular & Eintrag bearbeiten. \\
\bottomrule
\end{referenz}

\fussy
